\section{ÍNDICE TENTATIVO}

Estructura tentativa que organizará el desarrollo del documento final de tesis:

1. MARCO INTRODUCTORIO

1.1. Introducción y contexto cuántico

1.2. Antecedentes de la investigación

1.3. Planteamiento del problema

1.4. Objetivos de la investigación

1.5. Hipótesis y operacionalización de variables

1.6. Justificación (económica, social y científica)

1.7. Alcances y límites

1.8. Aportes esperados (práctico y teórico)

2. MARCO TEÓRICO

2.1. Criptografía basada en retículos (Lattices) y el problema LWE/RLWE

2.2. Estándares post-cuánticos del NIST: FIPS 203 (ML-KEM), FIPS 204 (ML-DSA) y FIPS 205 (SLH-DSA)

2.3. Pruebas de conocimiento cero: fundamentos y zk-SNARKs

2.4. El esquema Groth16 y la arithmetización R1CS

2.5. La función hash Poseidon en circuitos ZK-friendly

2.6. Blockchain, contratos inteligentes y la Máquina Virtual de Ethereum (EVM)

2.7. Modelos de votación electrónica con verificabilidad de extremo a extremo (E2E)

2.8. Integración híbrida: PQC, ZKP y Blockchain aplicados a la votación electrónica

3. MARCO APLICATIVO

3.1. Arquitectura general del sistema propuesto

3.2. Diseño del módulo criptográfico post-cuántico (registro, cifrado ML-KEM-768 y firma ML-DSA-65)

3.3. Diseño e implementación del circuito de conocimiento cero en Circom

3.4. Proceso de Trusted Setup y generación de pruebas Groth16 con snarkjs

3.5. Desarrollo de los contratos inteligentes: Bulletin Board (VotingE2E.sol) y Verifier.sol

3.6. Integración del pipeline completo de votación E2E

3.7. Configuración del entorno de simulación blockchain (Hardhat)

3.8. Desarrollo de la interfaz de línea de comandos (CLI) del votante

4. PRUEBA DE HIPÓTESIS

4.1. Diseño experimental y metodología de evaluación

4.2. Benchmarks de latencia del pipeline E2E

4.3. Medición de uso de CPU y memoria RAM

4.4. Pruebas de escalabilidad con hasta 10,000 votantes concurrentes

4.5. Análisis del costo de gas en la EVM

4.6. Contrastación de la hipótesis frente a los resultados obtenidos

4.7. Discusión de resultados y comparación con trabajos relacionados

5. CONCLUSIONES Y RECOMENDACIONES

5.1. Conclusiones generales

5.2. Verificación del cumplimiento de los objetivos específicos

5.3. Recomendaciones

5.4. Líneas de trabajo futuro
